\documentclass[10pt,letterpaper]{article}
\usepackage[utf8]{inputenc}
\usepackage[spanish]{babel}
\usepackage{amsmath,amssymb}
\usepackage{geometry}
\usepackage{array}
\usepackage[table]{xcolor}
\usepackage{booktabs}

\geometry{top=1.5cm, bottom=1.5cm, left=1.8cm, right=1.8cm}
\pagestyle{empty}


\definecolor{tableheader}{RGB}{30, 65, 110}   % Azul elegante
\definecolor{rowgray}{RGB}{245, 247, 250}     % Gris muy tenue

\begin{document}
	
	\begin{center}
		{\Large \textbf{Formulario: Ecuaciones Diferenciales Lineales de Orden Superior}}
	\end{center}
	
	\vspace{0.3cm}
	
	\section*{1. Solución General}
	\[
	y = y_c + y_p
	\]
	\begin{itemize}
		\item $y_c$: Solución homogénea ($R(x) = 0$).
		\item $y_p$: Solución particular debida a $R(x)$.
	\end{itemize}
	
	\hrulefill
	
	\section*{2. Solución Complementaria $y_c$}
	
	Se resuelve la ecuación auxiliar o característica sustituyendo las derivadas por potencias de $m$ ($y'' \to m^2$, $y' \to m$, $y \to 1$) e igualando a cero ($a_n m^n + \dots + a_1 m + a_0 = 0$).
	
	\begin{table}[!ht]
		\centering
		\renewcommand{\arraystretch}{1.5}
		\setlength{\tabcolsep}{10pt}
		\begin{tabular}{r >{$}l<{$} >{$}l<{$}}
			\rowcolor{tableheader}
			\multicolumn{1}{c}{\bfseries\color{white}\#} & 
			\multicolumn{1}{l}{\bfseries\color{white} Naturaleza de las raíces ($m$)} & 
			\multicolumn{1}{l}{\bfseries\color{white} Forma de la solución complementaria $y_c$} \\
			\midrule
			1 & m_1 \neq m_2 \text{ (reales distintas)} & C_1 e^{m_1 x} + C_2 e^{m_2 x} \\
			\rowcolor{rowgray}
			2 & m_1 \neq \dots \neq m_n \text{ (orden } n \text{ distintas)} & C_1 e^{m_1 x} + C_2 e^{m_2 x} + \dots + C_n e^{m_n x} \\
			3 & m_1 = m_2 = m \text{ (real repetida 2 veces)} & C_1 e^{mx} + C_2 x e^{mx} \\
			\rowcolor{rowgray}
			4 & m \text{ repetida } k \text{ veces (multiplicidad } k\text{)} & (C_1 + C_2 x + C_3 x^2 + \dots + C_k x^{k-1})e^{mx} \\
			5 & m = \pm \beta i \text{ (imaginarias puras, } \alpha = 0\text{)} & C_1 \cos(\beta x) + C_2 \sin(\beta x) \\
			\rowcolor{rowgray}
			6 & m = \alpha \pm \beta i \text{ (complejas conjugadas)} & e^{\alpha x}\big(C_1 \cos(\beta x) + C_2 \sin(\beta x)\big) \\
			7 & m = \alpha \pm \beta i \text{ (repetidas } k \text{ veces)} & e^{\alpha x}\big[(C_1 + \dots + C_k x^{k-1})\cos(\beta x) + (C_{k+1} + \dots + C_{2k} x^{k-1})\sin(\beta x)\big] \\
			\bottomrule
		\end{tabular}
		\caption{Formas de la solución complementaria $y_c$ según las raíces de la ecuación auxiliar.}
		\label{tab:solucion_complementaria}
	\end{table}
	
	\vspace{0.1cm}
	\noindent\textbf{Identificación para Variación de Parámetros:} Para orden 2, cada término que acompaña a una constante arbitraria define una función del conjunto fundamental:
	\[
	y_c = C_1 \mathbf{y_1} + C_2 \mathbf{y_2}
	\]
	
	
	\newpage
	
	\section*{3. Solución Particular $y_p$ (Coeficientes Indeterminados)}
	
	Para proponer la forma de $y_p$ a partir de la función $R(x)$, se definen las siguientes variables:
	\begin{itemize}
		\setlength{\itemsep}{1pt}
		\setlength{\parskip}{0pt}
		\item $k_0, k_1, k_2, \dots$: Números o coeficientes que acompañan a las $x$ en el polinomio de $R(x)$ ($k_0$ es el número solo).
		\item $\alpha$: Número que multiplica a la $x$ dentro del exponente ($e^{\alpha x}$).
		\item $\beta$: Número dentro del argumento del seno o coseno ($\sin \beta x$, $\cos \beta x$).
		\item $C_n, C_{n+1}, \dots$: Constantes por encontrar para armar $y_p$ (el índice $n$ continúa después de las $C$ usadas en $y_c$).
	\end{itemize}
	
	\begin{table}[!ht]
		\centering
		\renewcommand{\arraystretch}{1.45}
		\setlength{\tabcolsep}{6pt}
		\begin{tabular}{r >{$}l<{$} >{$}l<{$}}
			\rowcolor{tableheader}
			\multicolumn{1}{c}{\bfseries\color{white}\#} & 
			\multicolumn{1}{l}{\bfseries\color{white} Término no homogéneo $R(x)$} & 
			\multicolumn{1}{l}{\bfseries\color{white} Forma propuesta para $y_p$} \\
			\midrule
			1 & k \text{ (constante)} & C_n \\
			\rowcolor{rowgray}
			2 & k_1 x + k_0 & C_n x + C_{n+1} \\
			3 & k_2 x^2 + k_1 x + k_0 & C_n x^2 + C_{n+1} x + C_{n+2} \\
			\rowcolor{rowgray}
			4 & k_3 x^3 + k_2 x^2 + k_1 x + k_0 & C_n x^3 + C_{n+1} x^2 + C_{n+2} x + C_{n+3} \\
			5 & \sin(\beta x) & C_n \cos(\beta x) + C_{n+1} \sin(\beta x) \\
			\rowcolor{rowgray}
			6 & \cos(\beta x) & C_n \cos(\beta x) + C_{n+1} \sin(\beta x) \\
			7 & e^{\alpha x} & C_n e^{\alpha x} \\
			\rowcolor{rowgray}
			8 & (k_1 x + k_0)e^{\alpha x} & (C_n x + C_{n+1})e^{\alpha x} \\
			9 & (k_2 x^2 + k_1 x + k_0)e^{\alpha x} & (C_n x^2 + C_{n+1} x + C_{n+2})e^{\alpha x} \\
			\rowcolor{rowgray}
			10 & e^{\alpha x} \sin(\beta x) & e^{\alpha x}(C_n \cos(\beta x) + C_{n+1} \sin(\beta x)) \\
			11 & (k_2 x^2 + k_1 x + k_0)\sin(\beta x) & (C_n x^2 + C_{n+1} x + C_{n+2})\cos(\beta x) + (C_{n+3} x^2 + C_{n+4} x + C_{n+5})\sin(\beta x) \\
			\rowcolor{rowgray}
			12 & (k_1 x + k_0) e^{\alpha x} \cos(\beta x) & e^{\alpha x}\big[(C_n x + C_{n+1})\cos(\beta x) + (C_{n+2} x + C_{n+3})\sin(\beta x)\big] \\
			\bottomrule
		\end{tabular}
		\caption{Formas de la solución particular $y_p$ para coeficientes indeterminados.}
		\label{tab:coeficientes_indeterminados}
	\end{table}
	
	\vspace{0.2cm}
	\noindent\textbf{Regla de duplicidad:} Si algún término de $y_p$ ya aparece en la solución complementaria $y_c$, multiplica toda la familia de ese término por $x$ (o $x^2$, $x^s$) hasta eliminar la repetición.\\
	
	
	\textit{Ejemplo:} Si en $y_c = e^{\alpha x}(C_1\cos\beta x + C_2\sin\beta x)$ y $R(x) = e^{\alpha x}\cos\beta x$:  
	$\implies$ La propuesta estándar falla por duplicidad, por lo que se ajusta a:  
	$y_p = x e^{\alpha x}\big(C_n\cos\beta x + C_{n+1}\sin\beta x\big)$.
	
	\newpage
	
	\section*{4. Variacion de parametros}
	\subsection*{Paso 1: Forma estándar}
	Dada la ecuación lineal de segundo orden:
	\[
	a_2(x)y'' + a_1(x)y' + a_0(x)y = g(x)
	\]
	Se divide entre $a_2(x)$ para normalizar el coeficiente principal:
	\[
	y'' + P(x)y' + Q(x)y = R(x) \quad \text{donde} \quad R(x) = \frac{g(x)}{a_2(x)}
	\]
	
	\subsection*{Paso 2: Ecuación homogénea y conjunto fundamental}
	Para coeficientes constantes, se resuelve la ecuación auxiliar asociada en $m$:
	\[
	a_2 m^2 + a_1 m + a_0 = 0
	\]
	Se obtiene la solución complementaria $y_c = c_1 y_1 + c_2 y_2$ según la naturaleza de las raíces:
	\begin{itemize}
		\item \textbf{Raíces reales distintas ($m_1 \neq m_2$):}
		\[
		y_c = c_1 e^{m_1 x} + c_2 e^{m_2 x} \implies y_1 = e^{m_1 x}, \quad y_2 = e^{m_2 x}
		\]
		\item \textbf{Raíz real repetida ($m_1 = m_2 = m$):}
		\[
		y_c = c_1 e^{mx} + c_2 x e^{mx} \implies y_1 = e^{mx}, \quad y_2 = x e^{mx}
		\]
		\item \textbf{Raíces complejas conjugadas ($m = \alpha \pm \beta i$):}
		\[
		y_c = c_1 e^{\alpha x}\cos(\beta x) + c_2 e^{\alpha x}\sin(\beta x) \implies y_1 = e^{\alpha x}\cos(\beta x), \quad y_2 = e^{\alpha x}\sin(\beta x)
		\]
	\end{itemize}
	
	\subsection*{Paso 3: Cálculo del Wronskiano}
	A partir de $y_1$ y $y_2$, se evalúa el determinante:
	\[
	W(y_1, y_2) = \begin{vmatrix} y_1 & y_2 \\ y_1' & y_2' \end{vmatrix} = y_1 y_2' - y_1' y_2 \quad (W \neq 0)
	\]
	
	\subsection*{Paso 4: Determinación de $u_1$ y $u_2$}
	Aplicación directa de las fórmulas integrales (sin constantes de integración):
	\[
	u_1 = -\int \frac{y_2 \cdot R(x)}{W} \, dx
	\]
	\[
	u_2 = \int \frac{y_1 \cdot R(x)}{W} \, dx
	\]
	
	\subsection*{Paso 5: Solución particular y general}
	La solución particular se define como:
	\[
	y_p = u_1 y_1 + u_2 y_2
	\]
	La solución general completa del sistema es:
	\[
	y = y_c + y_p = c_1 y_1 + c_2 y_2 + u_1 y_1 + u_2 y_2
	\]
	
\end{document}